%==================================================
% MH1101 Chapter 1: Integrals
%==================================================

\section{Integrals}

\subsection{Antiderivatives and Indefinite Integrals}

\begin{definition}{Antiderivative and indefinite integral}{def:antiderivative}
Let $f$ be defined on an interval $I$. A function $F$ is an \textbf{antiderivative} of $f$ on $I$ if
\[
F'(x)=f(x)\quad \text{for all }x\in I.
\]
The \textbf{indefinite integral} denotes the family of all antiderivatives:
\[
\int f(x)\,dx = F(x)+C.
\]
\end{definition}

\begin{theorem}{All antiderivatives differ by a constant}{thm:antiderivatives-constant}
If $F$ and $G$ are antiderivatives of $f$ on an interval $I$, then $F(x)-G(x)$ is constant on $I$.
\end{theorem}

\begin{proof}
Since $F'(x)=f(x)=G'(x)$ on $I$, we have $(F-G)'(x)=0$ on $I$. Hence $F-G$ is constant on $I$.
\end{proof}

\begin{keyformula}[title={Key Formula: Basic antiderivative library}]
For constants $a\neq -1$ and $k\neq 0$:
\[
\int x^{a}\,dx=\frac{x^{a+1}}{a+1}+C,\qquad
\int e^{kx}\,dx=\frac{1}{k}e^{kx}+C,\qquad
\int \cos(kx)\,dx=\frac{1}{k}\sin(kx)+C,\qquad
\int \frac{1}{x}\,dx=\ln|x|+C.
\]
\end{keyformula}

\begin{examplebox}[title={Worked Example: General antiderivative (power simplification)}]
Find the most general antiderivative of $\displaystyle \frac{5-4x^3+2x^6}{x^6}$ for $x\neq 0$.

\textbf{Solution.}
Rewrite using powers:
\[
\frac{5-4x^3+2x^6}{x^6}=5x^{-6}-4x^{-3}+2.
\]
Integrate term-by-term:
\[
\int 5x^{-6}\,dx = 5\cdot \frac{x^{-5}}{-5}=-x^{-5},\qquad
\int (-4x^{-3})\,dx=-4\cdot \frac{x^{-2}}{-2}=2x^{-2},\qquad
\int 2\,dx=2x.
\]
Hence
\[
\int \frac{5-4x^3+2x^6}{x^6}\,dx = -\frac{1}{x^5}+\frac{2}{x^2}+2x+C.
\]
\end{examplebox}

\begin{commonmistake}[title={Common Mistake: Dropping $+C$ (and losing solutions)}]
Indefinite integrals represent \emph{families} of functions. Omitting $+C$ can remove valid solutions, especially when applying an initial condition later.
\end{commonmistake}

\begin{examtip}[title={Exam Focus: Solving for $f$ given $f'$ (and initial data)}]
Workflow:
\begin{enumerate}[label=(\arabic*)]
\item Integrate: $f(x)=\int f'(x)\,dx = F(x)+C$.
\item Apply the condition (e.g.\ $f(x_0)=y_0$) to solve for $C$.
\item If $f''$ is given, integrate twice: $f'(x)=\int f''(x)\,dx + C_1$, then $f(x)=\int f'(x)\,dx + C_2$.
\end{enumerate}
\end{examtip}

\subsection{Area Problem and Riemann Sums}

A key motivation for definite integrals is \textbf{area/accumulation}. For $f(x)\ge 0$ on $[a,b]$, the area under $y=f(x)$ from $a$ to $b$ can be approximated by rectangles.

Let $[a,b]$ be partitioned into $n$ equal subintervals with
\[
\Delta x = \frac{b-a}{n},\qquad x_i=a+i\Delta x\quad (i=0,1,\dots,n).
\]
Choosing sample points $x_i^\ast\in[x_{i-1},x_i]$, the associated \textbf{Riemann sum} is
\[
S_n=\sum_{i=1}^n f(x_i^\ast)\,\Delta x.
\]

% Figure description (fallback): A schematic curve y=f(x) above the x-axis on [a,b], with 4 right-endpoint rectangles whose heights are f(x_i). This illustrates a Riemann-sum approximation of area.

\begin{center}
\begin{tikzpicture}[scale=0.95]
\draw[->] (0,0) -- (6,0) node[right] {$x$};
\draw[->] (0,0) -- (0,3.4) node[above] {$y$};
% Interval [a,b] mapped to x in [1,5]
\draw (1,0) node[below] {$a$};
\draw (5,0) node[below] {$b$};
% A smooth schematic curve
\draw[thick,domain=1:5,samples=120,smooth]
plot (\x,{0.22*(\x-1)*(5-\x)+0.55});
% Right-endpoint rectangles for n=4, Delta x = 1 on [1,5]
\foreach \k in {1,2,3,4}{
\pgfmathsetmacro{\xl}{1+(\k-1)}
\pgfmathsetmacro{\xr}{1+\k}
\pgfmathsetmacro{\h}{0.22*(\xr-1)*(5-\xr)+0.55}
\draw[thin] (\xl,0) rectangle (\xr,\h);
}
\end{tikzpicture}
\end{center}

\subsection{Definite Integrals and Basic Properties}

\begin{definition}{Definite integral}{def:definite-integral}
Suppose $f$ is defined on $[a,b]$. Using the equal-width partition with $\Delta x=(b-a)/n$ and sample points $x_i^\ast\in[x_{i-1},x_i]$, the \textbf{definite integral} of $f$ on $[a,b]$ is
\[
\int_a^b f(x)\,dx := \lim_{n\to\infty}\sum_{i=1}^n f(x_i^\ast)\,\Delta x,
\]
provided the limit exists and is independent of the choice of sample points. In that case, $f$ is \textbf{integrable} on $[a,b]$.
\end{definition}

If $f$ changes sign, $\int_a^b f(x)\,dx$ is interpreted as \textbf{net area}: (area above $x$-axis) minus (area below $x$-axis).

\begin{keyformula}[title={Key Formula: Core properties of definite integrals}]
Assume $f,g$ are integrable on $[a,b]$ and $c$ is constant.
\begin{enumerate}[label=(P\arabic*), leftmargin=*]
\item $\displaystyle \int_a^b c\,dx = c(b-a)$.
\item $\displaystyle \int_a^b (f(x)\pm g(x))\,dx = \int_a^b f(x)\,dx \pm \int_a^b g(x)\,dx$.
\item $\displaystyle \int_a^b c f(x)\,dx = c\int_a^b f(x)\,dx$.
\item (Additivity) If $a<c<b$, then $\displaystyle \int_a^b f(x)\,dx = \int_a^c f(x)\,dx + \int_c^b f(x)\,dx$.
\item (Comparison) If $f(x)\ge 0$ on $[a,b]$, then $\displaystyle \int_a^b f(x)\,dx \ge 0$.
\item (Comparison) If $f(x)\ge g(x)$ on $[a,b]$, then $\displaystyle \int_a^b f(x)\,dx \ge \int_a^b g(x)\,dx$.
\item (Bounds) If $m\le f(x)\le M$ on $[a,b]$, then $\displaystyle m(b-a)\le \int_a^b f(x)\,dx \le M(b-a)$.
\end{enumerate}
\end{keyformula}

\begin{examplebox}[title={Worked Example: Average value satisfies its own bounds}]
Show that if $f$ is continuous on $[a,b]$ with $m\le f(x)\le M$, then the average value $\displaystyle f_{\mathrm{ave}}=\frac{1}{b-a}\int_a^b f(x)\,dx$ satisfies $m\le f_{\mathrm{ave}}\le M$.

\textbf{Solution.}
By Property P7,
\[
m(b-a)\le \int_a^b f(x)\,dx \le M(b-a).
\]
Divide by $b-a>0$ to obtain
\[
m \le \frac{1}{b-a}\int_a^b f(x)\,dx \le M,
\]
i.e.\ $m\le f_{\mathrm{ave}}\le M$.
\end{examplebox}

\begin{examtip}[title={Exam Focus: Fast integral comparisons}]
If you can compare integrands pointwise (e.g.\ $f\ge g\ge 0$), then you can compare integrals \emph{without computing antiderivatives}. This is often faster and avoids algebra mistakes.
\end{examtip}

\subsection{The Fundamental Theorem of Calculus}

\textbf{Motivation: Why do we need the FTC?}

At this point in the course, you know two distinct operations:
\begin{itemize}[leftmargin=*]
\item \textbf{Differentiation:} Given a function $f$, compute its instantaneous rate of change $f'(x)$.
\item \textbf{Integration (as area):} Given a function $f$ on $[a,b]$, compute the accumulated area $\int_a^b f(x)\,dx$ by taking the limit of Riemann sums.
\end{itemize}

These look completely unrelated---one is about \emph{rates}, the other is about \emph{totals}. Computing $\int_0^1 x^2\,dx$ from the definition means evaluating $\displaystyle \lim_{n\to\infty}\sum_{i=1}^n \left(\frac{i}{n}\right)^2\cdot\frac{1}{n}$, which requires summation formulas and careful limit-taking.

The Fundamental Theorem of Calculus reveals that differentiation and integration are \textbf{inverse processes}. This has two immediate practical consequences:
\begin{enumerate}[label=(\roman*), leftmargin=*]
\item If you \emph{integrate then differentiate}, you recover the original function (FTC Part~I).
\item If you know an \emph{antiderivative}, you can compute a definite integral instantly without Riemann sums (FTC Part~II).
\end{enumerate}

In effect, the FTC transforms integration from a \emph{limit-of-sums calculation} (slow, algebraically tedious) into an \emph{antiderivative evaluation} (fast, routine). This is why we can evaluate $\int_0^1 x^2\,dx$ by writing $\left[\tfrac{x^3}{3}\right]_0^1 = \tfrac{1}{3}$ instead of summing rectangles.

\vspace{0.5em}

We now state and prove both parts of the theorem carefully, emphasizing \emph{why} each formula takes the form it does.

\begin{theorem}{Fundamental Theorem of Calculus (Part I)}{thm:ftc1}
If $f$ is continuous on $[a,b]$ and
\[
F(x)=\int_a^x f(t)\,dt \qquad (a\le x\le b),
\]
then $F$ is differentiable on $(a,b)$ and
\[
F'(x)=f(x)\qquad (a<x<b).
\]
\end{theorem}

\textbf{Intuition behind Part I: Why does the derivative of an integral give back the original function?}

Think of $F(x)=\int_a^x f(t)\,dt$ as the \textbf{accumulated area} under $y=f(t)$ from $a$ up to $x$. When you increase $x$ by a small amount $h$, the new area is approximately the area of a thin rectangle of width $h$ and height $f(x)$, so
\[
F(x+h)-F(x) \approx f(x)\cdot h.
\]
Dividing by $h$ and taking $h\to 0$ gives $F'(x)=f(x)$. The proof below makes this intuition rigorous using the Mean Value Theorem for Integrals.

\begin{proof}
\textbf{Strategy:} Use the definition of the derivative and then apply the Mean Value Theorem for Integrals to evaluate the resulting integral increment.

\textbf{Step 1: Expand $F'(x)$ using the definition.}

Let $x\in(a,b)$. By definition,
\[
F'(x)=\lim_{h\to 0}\frac{F(x+h)-F(x)}{h}.
\]
Since $F(x)=\int_a^x f(t)\,dt$ and $F(x+h)=\int_a^{x+h} f(t)\,dt$, we have
\[
F(x+h)-F(x) = \int_a^{x+h} f(t)\,dt - \int_a^x f(t)\,dt.
\]
By the additivity property (P4), the interval $[a,x+h]$ splits as $[a,x]\cup[x,x+h]$, giving
\[
F(x+h)-F(x)=\int_x^{x+h} f(t)\,dt.
\]
Therefore
\[
F'(x)=\lim_{h\to 0}\frac{1}{h}\int_x^{x+h} f(t)\,dt.
\]

\textbf{Step 2: Apply the Mean Value Theorem for Integrals.}

Since $f$ is continuous on $[x,x+h]$ (for small $h>0$), the Mean Value Theorem for Integrals guarantees a point $c\in[x,x+h]$ such that
\[
\int_x^{x+h} f(t)\,dt = f(c)\cdot h.
\]
Substituting into our expression for $F'(x)$,
\[
F'(x)=\lim_{h\to 0}\frac{f(c)\cdot h}{h}=\lim_{h\to 0}f(c).
\]

\textbf{Step 3: Take the limit as $h\to 0$.}

Because $c\in[x,x+h]$, as $h\to 0$ we have $c\to x$ (the point $c$ is squeezed between $x$ and $x+h$). Since $f$ is continuous at $x$,
\[
\lim_{h\to 0}f(c)=f(x).
\]
Hence $F'(x)=f(x)$, as desired.
\end{proof}

\textbf{Why is continuity of $f$ required?}

The proof relies on two facts:
\begin{itemize}[leftmargin=*]
\item The Mean Value Theorem for Integrals, which requires $f$ to be continuous.
\item The limit $\lim_{c\to x}f(c)=f(x)$, which is the definition of continuity at $x$.
\end{itemize}
If $f$ has a jump discontinuity, the accumulated area function $F(x)$ may still be continuous, but it will fail to be differentiable at the jump.

\vspace{0.5em}

\textbf{Summary of Part I in words:} If $F(x)$ accumulates the values of $f$ from $a$ to $x$, then the \emph{rate of change} of this accumulation at $x$ is precisely $f(x)$ itself.

\begin{examplebox}[title={Worked Example: FTC Part I (direct application)}]
Find $\displaystyle \frac{d}{dx}\int_1^x \frac{1}{1+t^2}\,dt$.

\textbf{Solution.}
The integrand $f(t)=\frac{1}{1+t^2}$ is continuous on $\mathbb{R}$. By FTC Part I, with $a=1$,
\[
\frac{d}{dx}\int_1^x \frac{1}{1+t^2}\,dt = f(x)=\frac{1}{1+x^2}.
\]
\end{examplebox}

\begin{examplebox}[title={Worked Example: FTC Part I with Chain Rule}]
Find $\displaystyle \frac{d}{dx}\int_1^{x^4} \sec t\,dt$.

\textbf{Solution.}
Let $u=x^4$. Then
\[
\frac{d}{dx}\int_1^{x^4} \sec t\,dt = \frac{d}{dx}\int_1^u \sec t\,dt = \frac{d}{du}\left(\int_1^u \sec t\,dt\right)\cdot \frac{du}{dx}.
\]
By FTC Part I, $\displaystyle \frac{d}{du}\int_1^u \sec t\,dt = \sec u$. Since $\frac{du}{dx}=4x^3$, we obtain
\[
\frac{d}{dx}\int_1^{x^4} \sec t\,dt = \sec(x^4)\cdot 4x^3 = 4x^3\sec(x^4).
\]
\end{examplebox}

\begin{examplebox}[title={Worked Example: FTC Part I with two variable limits (Leibniz rule)}]
Find $\displaystyle \frac{d}{dx}\int_{x}^{x^2} \sin t\,dt$.

\textbf{Solution.}
\textbf{Method:} Split the integral at a fixed point (say, $0$) and apply FTC Part I to each piece.

Write
\[
\int_x^{x^2} \sin t\,dt = \int_x^0 \sin t\,dt + \int_0^{x^2} \sin t\,dt = -\int_0^x \sin t\,dt + \int_0^{x^2} \sin t\,dt.
\]
Differentiate term-by-term using FTC Part I and the Chain Rule:
\[
\frac{d}{dx}\left(-\int_0^x \sin t\,dt\right)=-\sin x,
\]
\[
\frac{d}{dx}\int_0^{x^2} \sin t\,dt = \sin(x^2)\cdot 2x.
\]
Hence
\[
\frac{d}{dx}\int_x^{x^2} \sin t\,dt = -\sin x + 2x\sin(x^2).
\]

\textbf{General form (Leibniz rule):}
\[
\frac{d}{dx}\int_{a(x)}^{b(x)} f(t)\,dt = f(b(x))\cdot b'(x) - f(a(x))\cdot a'(x).
\]
\end{examplebox}

\begin{examtip}[title={Exam Focus: Variable limits and Chain Rule}]
When the upper limit is not just $x$ but a composite function $g(x)$, you \emph{must} apply the Chain Rule:
\[
\frac{d}{dx}\int_a^{g(x)} f(t)\,dt = f(g(x))\cdot g'(x).
\]
When both limits are variable, split at a convenient constant and differentiate each piece separately.
\end{examtip}

\vspace{1em}

\begin{theorem}{Fundamental Theorem of Calculus (Part II)}{thm:ftc2}
If $f$ is continuous on $[a,b]$ and $F$ is \emph{any} antiderivative of $f$ (i.e., $F'(x)=f(x)$ on $[a,b]$), then
\[
\int_a^b f(x)\,dx = F(b)-F(a).
\]
\end{theorem}

\textbf{Intuition behind Part II: Why does an antiderivative let us evaluate definite integrals?}

Part I showed that the accumulation function $G(x)=\int_a^x f(t)\,dt$ is an antiderivative of $f$. If $F$ is \emph{any other} antiderivative of $f$, then $F$ and $G$ differ by a constant (Theorem~\ref{thm:antiderivatives-constant}):
\[
F(x)=G(x)+C.
\]
Evaluating at $x=b$,
\[
F(b)=G(b)+C=\int_a^b f(t)\,dt + C.
\]
Evaluating at $x=a$,
\[
F(a)=G(a)+C=\int_a^a f(t)\,dt + C = 0+C=C.
\]
Subtracting,
\[
F(b)-F(a)=\int_a^b f(t)\,dt.
\]
The constant $C$ cancels in the subtraction, which is why you can use \emph{any} antiderivative $F$ (with or without $+C$) to compute the definite integral.

\begin{proof}
\textbf{Strategy:} Define $G(x)=\int_a^x f(t)\,dt$. Both $G$ and $F$ are antiderivatives of $f$, so they differ by a constant. Exploit this to relate $\int_a^b f$ to $F(b)-F(a)$.

By FTC Part I, $G'(x)=f(x)$ on $(a,b)$. Since $F$ is also an antiderivative of $f$, we have $F'(x)=f(x)=G'(x)$ on $(a,b)$. By Theorem~\ref{thm:antiderivatives-constant}, there exists a constant $C$ such that
\[
F(x)=G(x)+C \qquad \text{for all }x\in[a,b].
\]

Evaluate at $x=b$:
\[
F(b)=G(b)+C=\int_a^b f(t)\,dt + C.
\]
Evaluate at $x=a$:
\[
F(a)=G(a)+C=\int_a^a f(t)\,dt + C = 0+C=C.
\]
Subtract the second equation from the first:
\[
F(b)-F(a)=\int_a^b f(t)\,dt.
\]
\end{proof}

\textbf{Why is the dummy variable irrelevant?}

In the statement of the theorem, we wrote $\int_a^b f(x)\,dx$ instead of $\int_a^b f(t)\,dt$. The choice of variable inside the integral doesn't matter:
\[
\int_a^b f(x)\,dx = \int_a^b f(t)\,dt = \int_a^b f(u)\,du.
\]
All represent the same number (the accumulated area from $a$ to $b$). The variable is a \textbf{dummy variable}---it is local to the integral and does not appear in the final answer $F(b)-F(a)$.

\vspace{0.5em}

\textbf{Notation:} We often write $F(b)-F(a)$ as $\left[F(x)\right]_a^b$ or $F(x)\Big|_a^b$.

\begin{examplebox}[title={Worked Example: FTC Part II (direct evaluation)}]
Evaluate $\displaystyle \int_{-1}^{2} x^3\,dx$.

\textbf{Solution.}
An antiderivative of $f(x)=x^3$ is $F(x)=\frac{x^4}{4}$ (we omit $+C$ because it will cancel). By FTC Part II,
\[
\int_{-1}^{2} x^3\,dx = \left[\frac{x^4}{4}\right]_{-1}^{2} = \frac{2^4}{4}-\frac{(-1)^4}{4}=\frac{16}{4}-\frac{1}{4}=\frac{15}{4}.
\]
\end{examplebox}

\begin{examplebox}[title={Worked Example: FTC Part II with endpoint checks}]
Evaluate $\displaystyle \int_0^\pi \cos x\,dx$.

\textbf{Solution.}
An antiderivative of $\cos x$ is $\sin x$. Hence
\[
\int_0^\pi \cos x\,dx = \left[\sin x\right]_0^\pi = \sin\pi - \sin 0 = 0-0=0.
\]

\textbf{Interpretation:} The area above the $x$-axis on $[0,\tfrac{\pi}{2}]$ exactly cancels the area below the $x$-axis on $[\tfrac{\pi}{2},\pi]$, giving net area zero.
\end{examplebox}

\begin{examplebox}[title={Worked Example: FTC Part II with absolute value (sign change)}]
Evaluate $\displaystyle \int_{-2}^2 |x|\,dx$.

\textbf{Solution.}
The function $|x|$ is not differentiable at $x=0$, so we must split the integral at the point where the formula changes:
\[
\int_{-2}^2 |x|\,dx = \int_{-2}^0 (-x)\,dx + \int_0^2 x\,dx.
\]
Compute each piece:
\[
\int_{-2}^0 (-x)\,dx = \left[-\frac{x^2}{2}\right]_{-2}^0 = 0-\left(-\frac{4}{2}\right)=2,
\]
\[
\int_0^2 x\,dx = \left[\frac{x^2}{2}\right]_0^2 = 2-0=2.
\]
Hence $\displaystyle \int_{-2}^2 |x|\,dx = 2+2=4$.

\textbf{Sanity check:} By symmetry, the two halves contribute equally; each is a triangle of base~2 and height~2, so area $\tfrac{1}{2}(2)(2)=2$.
\end{examplebox}

\begin{commonmistake}[title={Common Mistake: Ignoring sign changes in the integrand}]
If $f$ changes sign or formula on $[a,b]$, you \emph{cannot} apply FTC Part II blindly to $\int_a^b f(x)\,dx$. You must:
\begin{enumerate}[label=(\arabic*)]
\item Identify points $c$ where $f$ changes sign or formula.
\item Split the integral: $\int_a^b f = \int_a^c f + \int_c^b f$.
\item Evaluate each piece separately using the correct formula for $f$ on each subinterval.
\end{enumerate}
Example: $\int_{-1}^1 \frac{1}{x}\,dx$ is \textbf{not} $\ln|x|\big|_{-1}^1 = 0$. The integrand has a singularity at $x=0$; the integral is improper and divergent.
\end{commonmistake}

\begin{examtip}[title={Exam Focus: FTC Part II workflow}]
When evaluating $\int_a^b f(x)\,dx$:
\begin{enumerate}[label=(\arabic*), leftmargin=*]
\item Find \emph{any} antiderivative $F(x)$ of $f(x)$ (omit $+C$).
\item Compute $F(b)-F(a)$.
\item Sanity check:
\begin{itemize}[leftmargin=2em]
\item If $f\ge 0$ on $[a,b]$, the answer should be $\ge 0$.
\item If $f$ is odd and the interval is symmetric, the answer may be zero.
\item If $f$ changes sign, check that you've split correctly.
\end{itemize}
\end{enumerate}
\end{examtip}

\begin{examtip}[title={Exam Focus: When to use Part I vs.\ Part II}]
\textbf{Use FTC Part I when:}
\begin{itemize}[leftmargin=*]
\item You see $\displaystyle \frac{d}{dx}\int_a^{g(x)} f(t)\,dt$ (differentiate an integral).
\item The question asks for the derivative of an accumulation function.
\end{itemize}
\textbf{Use FTC Part II when:}
\begin{itemize}[leftmargin=*]
\item You need to evaluate $\int_a^b f(x)\,dx$ numerically.
\item You are asked to compute a definite integral using an antiderivative.
\end{itemize}
\end{examtip}

\subsection{The Substitution Rule}

Integration by substitution (also called $u$-substitution) is the reverse of the Chain Rule for differentiation. It allows us to evaluate integrals by transforming them into simpler integrals in a new variable.

\begin{theorem}{Substitution Rule (Indefinite Integrals)}{thm:substitution-indefinite}
If $u=g(x)$ is differentiable and $F$ is an antiderivative of $f$, then
\[
\int f(g(x))\,g'(x)\,dx = F(g(x))+C.
\]
Equivalently, writing $du=g'(x)\,dx$,
\[
\int f(u)\,du = F(u)+C.
\]
\end{theorem}

\textbf{Practical recipe:}
\begin{enumerate}[label=(\arabic*)]
\item Choose a substitution $u=g(x)$.
\item Compute $du=g'(x)\,dx$.
\item Rewrite the integrand and differential in terms of $u$ and $du$.
\item Integrate with respect to $u$.
\item Substitute back $u=g(x)$ to express the answer in terms of $x$.
\end{enumerate}

\begin{examplebox}[title={Worked Example: $u$-substitution for indefinite integrals}]
Evaluate $\displaystyle \int x\cos(x^2)\,dx$.

\textbf{Solution.}
Let $u=x^2$. Then $du=2x\,dx$, so $x\,dx=\frac{1}{2}du$. Substitute:
\[
\int x\cos(x^2)\,dx = \int \cos u\cdot \frac{1}{2}\,du = \frac{1}{2}\sin u + C = \frac{1}{2}\sin(x^2)+C.
\]
\end{examplebox}

\begin{theorem}{Substitution Rule (Definite Integrals)}{thm:substitution-definite}
If $g$ is continuously differentiable on $[a,b]$ and $f$ is continuous on the range of $g$, then
\[
\int_a^b f(g(x))\,g'(x)\,dx = \int_{g(a)}^{g(b)} f(u)\,du.
\]
\end{theorem}

\textbf{Key observation:} When you substitute $u=g(x)$ in a definite integral, you must also transform the limits of integration: $x=a$ becomes $u=g(a)$, and $x=b$ becomes $u=g(b)$.

\begin{examplebox}[title={Worked Example: $u$-substitution for definite integrals}]
Evaluate $\displaystyle \int_2^5 \frac{1}{3-5x}\,dx$.

\textbf{Solution.}
Let $u=3-5x$. Then $du=-5\,dx$, so $dx=-\frac{1}{5}du$.

Transform the limits:
\[
x=2 \Rightarrow u=3-5(2)=-7,\qquad x=5 \Rightarrow u=3-5(5)=-22.
\]
Substitute:
\[
\int_2^5 \frac{1}{3-5x}\,dx = \int_{-7}^{-22} \frac{1}{u}\cdot \left(-\frac{1}{5}\right)\,du = -\frac{1}{5}\int_{-7}^{-22} \frac{1}{u}\,du.
\]
Reverse the limits (introduces a minus sign):
\[
= \frac{1}{5}\int_{-22}^{-7} \frac{1}{u}\,du = \frac{1}{5}\left[\ln|u|\right]_{-22}^{-7} = \frac{1}{5}\bigl(\ln 7 - \ln 22\bigr)=\frac{1}{5}\ln\left(\frac{7}{22}\right).
\]
\end{examplebox}

\begin{examtip}[title={Exam Focus: Substitution workflow for definite integrals}]
When using substitution on $\int_a^b f(x)\,dx$:
\begin{enumerate}[label=(\arabic*), leftmargin=*]
\item Choose $u=g(x)$ and compute $du=g'(x)\,dx$.
\item Transform limits: $u_1=g(a)$, $u_2=g(b)$.
\item Rewrite the integral entirely in terms of $u$ (no $x$ should remain).
\item Integrate with respect to $u$ from $u_1$ to $u_2$.
\item \textbf{Do not} substitute back to $x$---evaluate directly in $u$.
\end{enumerate}
\end{examtip}

\newpage
\subsection{Improper Integrals}

So far, we have defined $\int_a^b f(x)\,dx$ only when $[a,b]$ is a finite interval and $f$ is continuous (or at least bounded and integrable) on $[a,b]$. An \textbf{improper integral} extends the concept of integration to:
\begin{itemize}
\item \textbf{Type 1:} Infinite intervals (e.g., $\int_1^\infty \frac{1}{x^2}\,dx$).
\item \textbf{Type 2:} Unbounded integrands (e.g., $\int_0^1 \frac{1}{\sqrt{x}}\,dx$, where $f(x)\to\infty$ as $x\to 0^+$).
\end{itemize}

\begin{definition}{Improper Integral of Type 1 (Infinite Interval)}{def:improper-type1}
\begin{enumerate}[label=(\alph*), leftmargin=*]
\item If $\int_a^t f(x)\,dx$ exists for all $t\ge a$, then
\[
\int_a^\infty f(x)\,dx := \lim_{t\to\infty}\int_a^t f(x)\,dx.
\]
\item If $\int_t^b f(x)\,dx$ exists for all $t\le b$, then
\[
\int_{-\infty}^b f(x)\,dx := \lim_{t\to-\infty}\int_t^b f(x)\,dx.
\]
\item If both $\int_{-\infty}^a f$ and $\int_a^\infty f$ converge, then
\[
\int_{-\infty}^\infty f(x)\,dx := \int_{-\infty}^a f(x)\,dx + \int_a^\infty f(x)\,dx.
\]
\end{enumerate}
In each case, the integral \textbf{converges} if the limit exists (and is finite); otherwise it \textbf{diverges}.
\end{definition}

\begin{keyformula}[title={Key Formula: $p$-integrals (Type 1 and Type 2)}]
\[
\int_1^\infty \frac{1}{x^p}\,dx \text{ converges } \iff p>1,
\qquad
\int_0^1 \frac{1}{x^p}\,dx \text{ converges } \iff p<1.
\]
\end{keyformula}

% Figure description (fallback): Plot of y=1/x^2 for x>=1, decreasing to 0; the "tail area" from 1 to infinity is finite, illustrating a convergent improper integral.

\begin{center}
\begin{tikzpicture}[scale=0.95]
\draw[->] (0.5,0) -- (6.2,0) node[right] {$x$};
\draw[->] (1,0) -- (1,3.3) node[above] {$y$};
\draw (1,0) node[below] {$1$};
\draw[thick,domain=1:6,samples=120,smooth]
plot (\x,{2/(\x*\x)});
\draw[dashed] (1,2) -- (1,0);
\node at (4.9,2.7) {$y=\frac{1}{x^2}$ (schematic)};
\end{tikzpicture}
\end{center}

\begin{examplebox}[title={Worked Example: Type 1 divergence}]
Determine whether $\displaystyle \int_1^\infty \frac{1}{x}\,dx$ converges.

\textbf{Solution.}
Write the definition:
\[
\int_1^\infty \frac{1}{x}\,dx=\lim_{t\to\infty}\int_1^t \frac{1}{x}\,dx
=\lim_{t\to\infty}\bigl[\ln|x|\bigr]_{1}^{t}
=\lim_{t\to\infty}\ln t=\infty.
\]
Hence the improper integral \textbf{diverges}.
\end{examplebox}

\begin{examplebox}[title={Worked Example: Type 1 evaluation (split correctly)}]
Evaluate $\displaystyle \int_{-\infty}^{\infty}\frac{1}{1+x^2}\,dx$.

\textbf{Solution.}
Split at $0$ and take one-sided limits:
\[
\int_{-\infty}^{\infty}\frac{1}{1+x^2}\,dx
=\int_{-\infty}^{0}\frac{1}{1+x^2}\,dx+\int_{0}^{\infty}\frac{1}{1+x^2}\,dx,
\]
provided both sides converge.

Since $\int \frac{1}{1+x^2}\,dx=\tan^{-1}x+C$,
\[
\int_{0}^{\infty}\frac{1}{1+x^2}\,dx=\lim_{t\to\infty}\bigl[\tan^{-1}x\bigr]_0^t
=\frac{\pi}{2},
\qquad
\int_{-\infty}^{0}\frac{1}{1+x^2}\,dx=\lim_{t\to-\infty}\bigl[\tan^{-1}x\bigr]_t^0
=\frac{\pi}{2}.
\]
Therefore $\displaystyle \int_{-\infty}^{\infty}\frac{1}{1+x^2}\,dx=\pi$.
\end{examplebox}

\begin{definition}{Improper Integral of Type 2 (Unbounded Integrand)}{def:improper-type2}
\begin{enumerate}[label=(\alph*), leftmargin=*]
\item If $f$ is continuous on $[a,b)$ and $\lim_{x\to b^-}f(x)=\pm\infty$, then
\[
\int_a^b f(x)\,dx := \lim_{t\uparrow b}\int_a^t f(x)\,dx.
\]
\item If $f$ is continuous on $(a,b]$ and $\lim_{x\to a^+}f(x)=\pm\infty$, then
\[
\int_a^b f(x)\,dx := \lim_{t\downarrow a}\int_t^b f(x)\,dx.
\]
\item If $f$ has a discontinuity at $c\in(a,b)$ and both $\int_a^c f$ and $\int_c^b f$ (as improper integrals) converge, then
\[
\int_a^b f(x)\,dx := \int_a^c f(x)\,dx + \int_c^b f(x)\,dx.
\]
\end{enumerate}
The integral converges if the limit(s) exist and are finite; otherwise it diverges.
\end{definition}

\begin{examplebox}[title={Worked Example: Type 2 at an endpoint}]
Evaluate $\displaystyle \int_{2}^{5}\frac{1}{\sqrt{x-2}}\,dx$.

\textbf{Solution.}
This is improper at $x=2$. By definition,
\[
\int_{2}^{5}\frac{1}{\sqrt{x-2}}\,dx=\lim_{t\downarrow 2}\int_{t}^{5}\frac{1}{\sqrt{x-2}}\,dx.
\]
Let $u=x-2$ so $du=dx$. Then $\int u^{-1/2}\,du = 2u^{1/2}$, hence
\[
\int_{t}^{5}\frac{1}{\sqrt{x-2}}\,dx
= \left[2\sqrt{x-2}\right]_{t}^{5}
=2\sqrt{3}-2\sqrt{t-2}.
\]
Taking $t\downarrow 2$ gives $2\sqrt{t-2}\to 0$, so the integral converges and equals $2\sqrt{3}$.
\end{examplebox}

\begin{commonmistake}[title={Common Mistake: Treating $\int_{-\infty}^{\infty}$ as a symmetric single limit}]
The definition is \emph{not} $\lim_{A\to\infty}\int_{-A}^{A} f(x)\,dx$.

A two-sided improper integral converges only if \textbf{both} one-sided integrals converge:
\[
\int_{-\infty}^{\infty} f := \int_{-\infty}^{a} f + \int_a^{\infty} f,
\quad \text{with both terms finite.}
\]
Symmetric cancellation can produce a finite number even when the improper integral diverges (see extension below).
\end{commonmistake}

\begin{examtip}[title={Exam Focus: Improper integrals (non-negotiable workflow)}]
\begin{enumerate}[label=(\arabic*), leftmargin=*]
\item Identify why it is improper (infinite interval / unbounded integrand / interior singularity).
\item Rewrite using the correct one-sided limit(s) \emph{before} manipulating.
\item Decide convergence; only then compute the finite value (if it converges).
\end{enumerate}
\end{examtip}

\begin{extension}[title={Extension (Tutorial 3, Question 2, Method 1)}]
\textbf{Syllabus link:} Chapter 1 — \emph{§1.6 Improper Integrals}.\\
\textbf{Proposition (independence of split point).}

Assume that for every real $c$, both $\int_{-\infty}^{c} f(x)\,dx$ and $\int_{c}^{\infty} f(x)\,dx$ converge.
Then for any real $a,b$,
\[
\int_{-\infty}^{a} f(x)\,dx+\int_{a}^{\infty} f(x)\,dx
=
\int_{-\infty}^{b} f(x)\,dx+\int_{b}^{\infty} f(x)\,dx.
\]

\textbf{Proof.} WLOG let $a<b$. For $M>0$, finite additivity gives
\[
\int_{-M}^{b} f=\int_{-M}^{a} f+\int_{a}^{b} f.
\]
Take $M\to\infty$ (the limits exist by hypothesis) to obtain
\[
\int_{-\infty}^{b} f=\int_{-\infty}^{a} f+\int_{a}^{b} f. \tag{1}
\]
Similarly, for $N>0$, finite additivity gives $\int_{a}^{N} f=\int_{a}^{b} f+\int_{b}^{N} f$; letting $N\to\infty$ yields
\[
\int_{a}^{\infty} f=\int_{a}^{b} f+\int_{b}^{\infty} f. \tag{2}
\]
Rearrange (2) as $\int_{b}^{\infty} f=\int_{a}^{\infty} f-\int_{a}^{b} f$ and add to (1):
\[
\int_{-\infty}^{b} f+\int_{b}^{\infty} f
=\left(\int_{-\infty}^{a} f+\int_{a}^{b} f\right)+\left(\int_{a}^{\infty} f-\int_{a}^{b} f\right)
=\int_{-\infty}^{a} f+\int_{a}^{\infty} f.
\]
\end{extension}

\begin{extension}[title={Extension (Beyond Syllabus: Cauchy principal value vs.\ improper integral)}]
\textbf{Syllabus link:} Chapter 1 — \emph{§1.6 Improper Integrals}.\\
\textbf{Definition.} For integrable $f$ on every finite interval, the \emph{Cauchy principal value} is
\[
\operatorname{PV}\!\int_{-\infty}^{\infty} f(x)\,dx := \lim_{A\to\infty}\int_{-A}^{A} f(x)\,dx
\quad\text{(if the limit exists).}
\]
This is \emph{not} the same as the improper integral, which requires both $\int_{-\infty}^{0} f$ and $\int_{0}^{\infty} f$ to converge (finite).\\
\textbf{Example (odd integrand with finite PV but divergent improper integral).}

Let $f(x)=\frac{x}{1+x^2}$ (odd, continuous on $\mathbb{R}$).
Compute one-sided integrals:
\[
\int_{0}^{A}\frac{x}{1+x^2}\,dx=\left[\frac12\ln(1+x^2)\right]_{0}^{A}=\frac12\ln(1+A^2)\xrightarrow[A\to\infty]{}\infty,
\]
\[
\int_{-A}^{0}\frac{x}{1+x^2}\,dx=\left[\frac12\ln(1+x^2)\right]_{-A}^{0}= -\frac12\ln(1+A^2)\xrightarrow[A\to\infty]{}-\infty.
\]
Hence the improper integral $\int_{-\infty}^{\infty} f$ \emph{does not converge} (one-sided integrals are not both finite).

But for every $A>0$, oddness gives $\int_{-A}^{A} f(x)\,dx=0$, so
\[
\operatorname{PV}\!\int_{-\infty}^{\infty}\frac{x}{1+x^2}\,dx=\lim_{A\to\infty}0=0.
\]
\end{extension}

%==================================================
% USER COMMENTS (EDITABLE)
%==================================================
% The FTC section has been enhanced with:
% - A motivational paragraph explaining *why* the FTC is needed (bridging rates and totals).
% - Intuition paragraphs *before* each part, explaining the geometric/physical reasoning.
% - Detailed proof annotations with step-by-step strategy and "why this works" commentary.
% - Explicit discussion of continuity hypothesis and dummy-variable conventions.
% - Additional worked examples with FTC Part I (Chain Rule, Leibniz rule) and Part II (absolute value, sign-change splitting).
% - Common mistakes and exam tips focusing on when to use Part I vs Part II, and how to handle variable limits.
% - All existing content preserved; only explanatory material and examples added.
%
% User request: "add explanation to the FTC part especially, avoid giving the definitions directly without any explanation on why are they like that"
% This has been addressed by:
%   1. Adding a "Motivation" subsection before the FTC theorems.
%   2. Including "Intuition behind Part I" and "Intuition behind Part II" paragraphs that explain the geometric/physical meaning *before* the formal statements.
%   3. Annotating the proofs with strategy comments and step-by-step reasoning.
%   4. Adding "Why is X required?" discussions after each theorem.
%   5. Including multiple worked examples demonstrating Chain Rule, Leibniz rule, absolute values, and sanity checks.
%==================================================
